\begin{ex}
Lít được viết tắt bằng chữ cái nào trong bảng chữ cái? Viết tắt có mấy chữ cái?
\shortans{1}
\loigiai{
Viết tắt là chữ l (chỉ gồm 1 chữ cái).
}
\end{ex}

\begin{ex}
Tính: $4\text{ }l + 5\text{ }l$. Kết quả là bao nhiêu lít?
\shortans{9}
\loigiai{
Ta có: $4 + 5 = 9$ lít.
}
\end{ex}

\begin{ex}
Tính: $12\text{ }l + 7\text{ }l$. Kết quả là bao nhiêu lít?
\shortans{19}
\loigiai{
Ta có: $12 + 7 = 19$ lít.
}
\end{ex}

\begin{ex}
Tính: $35\text{ }l + 20\text{ }l$. Kết quả là bao nhiêu lít?
\shortans{55}
\loigiai{
Ta có: $35 + 20 = 55$ lít.
}
\end{ex}

\begin{ex}
Tính: $18\text{ }l - 8\text{ }l$. Kết quả là bao nhiêu lít?
\shortans{10}
\loigiai{
Ta có: $18 - 8 = 10$ lít.
}
\end{ex}

\begin{ex}
Tính: $47\text{ }l - 23\text{ }l$. Kết quả là bao nhiêu lít?
\shortans{24}
\loigiai{
Ta có: $47 - 23 = 24$ lít.
}
\end{ex}

\begin{ex}
Tính: $50\text{ }l - 30\text{ }l$. Kết quả là bao nhiêu lít?
\shortans{20}
\loigiai{
Ta có: $50 - 30 = 20$ lít.
}
\end{ex}

\begin{ex}
Tính giá trị biểu thức: $15\text{ }l + 5\text{ }l - 8\text{ }l$.
\shortans{12}
\loigiai{
Ta có $15 + 5 = 20$, rồi $20 - 8 = 12$ lít.
}
\end{ex}

\begin{ex}
Tính giá trị biểu thức: $30\text{ }l - 10\text{ }l + 15\text{ }l$.
\shortans{35}
\loigiai{
Ta có $30 - 10 = 20$, rồi $20 + 15 = 35$ lít.
}
\end{ex}

\begin{ex}
Thùng thứ nhất có 16 lít nước ngọt, thùng thứ hai có 12 lít nước ngọt. Cả hai thùng có tất cả bao nhiêu lít nước ngọt?
\shortans{28}
\loigiai{
Cả hai thùng có: $16 + 12 = 28$ lít.
}
\end{ex}

\begin{ex}
Thùng thứ nhất nhiều hơn thùng thứ hai bao nhiêu lít nước ngọt?
\shortans{4}
\loigiai{
Nhiều hơn: $16 - 12 = 4$ lít.
}
\end{ex}

\begin{ex}
Can dầu to có 25 lít dầu, rót sang can nhỏ 10 lít dầu. Can dầu to còn lại bao nhiêu lít dầu?
\shortans{15}
\loigiai{
Còn lại: $25 - 10 = 15$ lít.
}
\end{ex}

\begin{ex}
Một bể nước có 80 lít nước, sau khi dùng để tưới cây thì bể còn 50 lít nước. Đã dùng bao nhiêu lít nước để tưới cây?
\shortans{30}
\loigiai{
Đã dùng: $80 - 50 = 30$ lít nước.
}
\end{ex}

\begin{ex}
Trong bình có 2 lít nước mơ, mẹ pha thêm 3 lít nước lọc vào bình. Trong bình có tất cả bao nhiêu lít nước?
\shortans{5}
\loigiai{
Có tất cả: $2 + 3 = 5$ lít.
}
\end{ex}

\begin{ex}
Bác bảo vệ mua 2 can xăng, mỗi can 5 lít. Hỏi bác mua tất cả bao nhiêu lít xăng?
\shortans{10}
\loigiai{
Bác mua: $5 + 5 = 10$ lít xăng.
}
\end{ex}

\begin{ex}
Một xô nước đựng được 10 lít nước, một chậu nước đựng được 15 lít nước. Cả xô và chậu đựng được bao nhiêu lít nước?
\shortans{25}
\loigiai{
Đựng được: $10 + 15 = 25$ lít.
}
\end{ex}

\begin{ex}
Điền số thích hợp: $14\text{ }l + \dots\text{ }l = 20\text{ }l$.
\shortans{6}
\loigiai{
Số cần điền là: $20 - 14 = 6$.
}
\end{ex}

\begin{ex}
Điền số thích hợp: $38\text{ }l - \dots\text{ }l = 25\text{ }l$.
\shortans{13}
\loigiai{
Số cần điền là: $38 - 25 = 13$.
}
\end{ex}

\begin{ex}
Điền số thích hợp: $\dots\text{ }l + 22\text{ }l = 45\text{ }l$.
\shortans{23}
\loigiai{
Số cần điền là: $45 - 22 = 23$.
}
\end{ex}

\begin{ex}
Một máy giặt lần giặt thứ nhất dùng hết 35 lít nước, lần giặt thứ hai dùng hết 30 lít nước. Cả hai lần giặt dùng hết bao nhiêu lít nước?
\shortans{65}
\loigiai{
Dùng hết: $35 + 30 = 65$ lít nước.
}
\end{ex}

\begin{ex}
Một thùng sữa chua có 24 lít sữa, đã bán được 14 lít. Thùng còn lại bao nhiêu lít sữa chua?
\shortans{10}
\loigiai{
Còn lại: $24 - 14 = 10$ lít.
}
\end{ex}

\begin{ex}
Có 4 chai nước khoáng, mỗi chai chứa 2 lít nước. Cả 4 chai chứa tất cả bao nhiêu lít nước?
\shortans{8}
\loigiai{
Chứa: $2 + 2 + 2 + 2 = 8$ lít.
}
\end{ex}

\begin{ex}
Một téc nước mini có 95 lít nước, đã xả bớt 42 lít để cọ rửa sân. Téc nước còn lại bao nhiêu lít?
\shortans{53}
\loigiai{
Còn lại: $95 - 42 = 53$ lít.
}
\end{ex}

\begin{ex}
Bình nước nhà bạn An có 18 lít, hôm nay dùng hết 6 lít. Bình còn lại bao nhiêu lít nước?
\shortans{12}
\loigiai{
Còn lại: $18 - 6 = 12$ lít.
}
\end{ex}

\begin{ex}
Tính: $10\text{ }l + 20\text{ }l + 30\text{ }l - 15\text{ }l$.
\shortans{45}
\loigiai{
Ta có $10 + 20 + 30 = 60$, rồi $60 - 15 = 45$ lít.
}
\end{ex}
